\section{Problem Statement}
%=====================================================================

This project addresses the divided access and facility-use process at the Gimnasio de Pesas. Authorization depends on separate sources (WooCommerce for paid memberships and administrative records for athletes) that staff must verify manually, and check-ins do not produce a centralized count of active users or consistent utilization metrics.

This is a relevant problem because the manual process is outdated and needed updating, and stakeholder discussions with UPRM Athletic Activities Director José F. Estévez confirmed the need and supported its selection as a Capstone project. It requires the design, implementation, and testing of an integrated software system. The software must determine authorization from multiple eligibility sources, enforce permissions for users, staff, and administrators, maintain active facility sessions, and transform recorded facility use into consistent utilization metrics. It also presents engineering challenges such as integrating with WooCommerce; keeping sessions and occupancy consistent under simultaneous requests; protecting sensitive user information; and handling imprecise indoor location and loss of connectivity. These behaviors can be evaluated through measurable software requirements and acceptance tests.


The target domain is university athletic-facility access and utilization operations. This domain covers the operational processes used to determine who is authorized to use a facility, register facility entry and use, maintain current operating schedules and temporary operational changes, and collect facility-use records that can support estimated occupancy and the identification of utilization patterns. This project belongs to that domain because it addresses how an athletic facility controls access, records facility use, communicates current operational information, and uses those records to support daily operations and administrative decision-making. For the current core scope, the project addresses these processes for the Gimnasio de Pesas located in the UPRM Natatorio Complex. Out of scope are payment processing, which remains in the existing WooCommerce process; access workflows for other facilities, including the pool, which have separate rules and authorization needs; medical-record management, since only safety information UPRM Athletic Activities deems necessary is kept; and capacity-based admission control, which requires an official UPRM capacity policy.

The intended users of the system fall into two groups. Facility users are authorized individuals who use the Gimnasio de Pesas through a paid membership or athlete authorization. Operational users include facility staff, including student workers assigned to reception or check-in duties, who monitor or assist with facility entry, and administrators who manage authorization records, athlete rosters, schedules, access rules, and utilization information. The principal stakeholders include the Departamento de Actividades Atléticas, because it owns the operational need being addressed; Prof. José F. Estévez, because he represents the client in project decisions; facility staff and administrators, because the system affects the processes they use to verify access and manage the facility; and authorized facility users, because the system affects how their access and facility use are recorded and communicated.

The proposed innovation is an integrated access and facility-use management platform for the Gimnasio de Pesas. It connects institutional login, access authorization (paid semester membership or preapproved athlete roster), location-validated check-in, active sessions, estimated occupancy, operating schedules, and utilization analytics in one workflow. It leaves the WooCommerce payment process in place and uses completed membership orders to determine paid access. A successful check-in creates a facility-use record that supports current operations and later analysis.

%=====================================================================
